\documentclass[11pt]{article}
\usepackage[margin=1in]{geometry}
\usepackage{amsmath,amssymb,amsthm}
\usepackage[T1]{fontenc}
\usepackage[hidelinks]{hyperref}
\newtheorem{theorem}{Theorem}
\newtheorem{lemma}{Lemma}
\newtheorem{remark}{Remark}
\newcommand{\N}{\mathbb N_0}
\newcommand{\LL}{\mathcal L}
\newcommand{\E}{\mathbb E}
\title{All-load stability of a deterministic-first-chunk file-sharing chain}
\author{Alec Kriebel\\Independent researcher\\
\small\url{https://orcid.org/0009-0001-9320-500X}}
\date{20 October 2026\\\small Unrefereed preprint draft}
\begin{document}
\maketitle
\begin{abstract}
We prove positive recurrence at every fixed positive arrival rate for the
four-population deterministic-first-chunk chain described by Norros, Reittu
and Eirola. This answers the stability conjecture in Norros's 2010
Oberwolfach contribution and the retained 2011 author proof, for the exact
six-rate model with a permanent seed and invisible empty waiting rooms.
A single coercive Foster function combines remaining workload, rounded
cohort imbalance, and two bounded finite birth--death Poisson correctors.
Monotone ratio cutoffs make every first-download interface error
nonpositive. The other interface errors are controlled pointwise.
The proof requires no stationary averaging, moment assumption, or
uniformity in the arrival rate. Exact finite controls supplement the
analytic proof. The contribution is the all-load extension of the stated historical
conjecture; classical methods and source-version limits are explicit.
\end{abstract}

\section{The model and the question}

An arriving peer chooses, independently and uniformly, which of two
chunks to download first. Let $a,b$ count empty peers assigned to the two
choices, and let $x,y$ count peers possessing just the respective chunks.
A completed peer departs immediately. There is one permanent seed.
The empty waiting rooms are invisible to the sampling denominator.
Precisely, for $\lambda>0$ and $D=x+y+1$, the chain on $\N^4$ has transitions
\[
\begin{array}{c|c}
\text{increment in }(a,b,x,y)&\text{rate}\\ \hline
(1,0,0,0)&\lambda/2\\
(0,1,0,0)&\lambda/2\\
(-1,0,1,0)&a(x+1)/D\\
(0,-1,0,1)&b(y+1)/D\\
(0,0,-1,0)&x(y+1)/D\\
(0,0,0,-1)&y(x+1)/D .
\end{array}                                                    \tag{1}
\]
A decrement at a zero coordinate has rate zero. These are the rates of
Figure 2 in \cite{NRE09} and the checked 2011 author proof of
\cite{NRE11}; the seed produces the $+1$ terms.

Norros \cite[pp.~2782--2783]{OWR} conjectured stochastic stability despite
instability of the corresponding large-system differential equations.
The retained January 5, 2011 author proof states the stability conjecture
in Section 3.2 as Conjecture 3.6. The earlier 2009 preprint
\cite[Conjecture 3.4]{NRE09} conjectured high-load instability instead.
These editions must be distinguished. Stability here means
irreducibility and positive recurrence in continuous time, as in those
sources.

\begin{theorem}\label{thm:main}
For every fixed real $\lambda>0$, the chain \textup{(1)} is nonexplosive,
irreducible and positive recurrent. It therefore has a unique stationary
probability distribution.
\end{theorem}

The exact rates are part of the theorem. In particular, no bound on the
seed's aggregate upload rate is imposed beyond \textup{(1)}. We assert no
uniform bound in $\lambda$, finite stationary population moment,
geometric ergodicity, or conclusion for other chunk-selection protocols.
The classical Foster method is background \cite{MT}. The work below is
the all-load pointwise certificate for this chain.

Write $\mathsf n=a+b+x+y$. A short observation already handles small
loads. For the workload $W=2a+2b+x+y$, direct calculation gives
\[
 \LL W
 =2\lambda-\frac{a(x+1)+b(y+1)+2xy+x+y}{D}
 \le 2\lambda-\frac{\mathsf n}{\mathsf n+1}.
                                                               \tag{2}
\]
Thus the classical drift criterion gives stability when
$0<\lambda<1/2$. The existence of some stable parameter values, even
with an unstable large-system limit, does not require our all-load
argument. Our aim is the extension to every positive load.

\section{Workload, imbalance, and a finite corrector}

For the rest of the proof fix $\lambda>0$ and put
\[
 Z=a+x-b-y,\quad
 u=\frac{a(x+1)+b(y+1)}D,\qquad
 d=\frac{2xy+x+y}D.
\]
First downloads leave $Z$ unchanged. Its other jumps have size one and
total rate $\lambda+d$. Hence
\[
 \LL W=2\lambda-u-d,\qquad \LL Z=\frac{y-x}D.                    \tag{3}
\]
For $m>0$ define the rounded absolute value
\[
 h_m(z)=
 \begin{cases}
 z^2/(2m)+m/2,&|z|\le m,\\
 |z|,&|z|\ge m.
 \end{cases}
\]
Its derivative is $z/m$ clipped to $[-1,1]$, and is globally
$1/m$-Lipschitz. Therefore, also across the rounding boundaries,
\[
 h_m(z+v)-h_m(z)\le h'_m(z)v+\frac{v^2}{2m}.
\]
Set
\[
 \beta=\frac{17\lambda}{8},\quad q=\frac78,\quad
 c=2(\beta+4),\quad m=4c,\quad \theta=\beta+c+4.
\]
Equations (3) give
\[
 \LL(W+ch_m(Z))
 \le \beta-u-qd+c h'_m(Z)\frac{y-x}D.                           \tag{4}
\]
The last term has absolute value at most $c$. If $Z\ge m$ and
$x\ge3(y+1)$, it is at most $-c/2$.

Put
\[
 \gamma=\frac{\theta+1}{\theta+2},\qquad
 \eta=\frac{1+\gamma}{2},\qquad \ell=\frac1{\eta-\gamma}.
\]
Consider on $\{0,\ldots,R\}$ the reflected birth--death generator $Q_R$
with births $\gamma(k+1)$ for $k<R$ and deaths $k$ for $k>0$.
Its stationary weights are $\gamma^k$ and its mean is
\[
 \mu_R=\frac{\sum_{k=0}^R k\gamma^k}{\sum_{k=0}^R\gamma^k}.
\]
The means increase to $\gamma/(1-\gamma)=\theta+1$. We can therefore
choose a finite integer $R\ge1$ such that
\[
 \mu:=\mu_R>\theta,\qquad q(R+1)>\beta+c+2.                     \tag{5}
\]
For $0\le k<R$ put
\[
 t_k=\frac{\sum_{j=0}^k\gamma^j(j-\mu)}
                  {\gamma^{k+1}(k+1)},\qquad
 g(k)=-\sum_{j=k}^{R-1}t_j,
\]
and let $g(k)=0$ for all integers $k\ge R$.

\begin{lemma}\label{lem:poisson}
The function $g$ is nonnegative and decreasing. With $G=g(0)>0$,
\[
 Q_Rg(k)=k-\mu\quad(0\le k\le R).
\]
For the actual chain,
\[
 \LL g(y)\le y-\mu
 \quad\text{if }0\le y\le R,\quad b/D\ge\gamma.                 \tag{6}
\]
For $y>R$ all increments of $g(y)$ vanish. The same assertions hold
with $(b,y)$ replaced by $(a,x)$.
\end{lemma}
\begin{proof}
The mean restricted to the prefix $\{0,\ldots,k\}$ is strictly below
$\mu$, so $t_k<0$. This proves monotonicity and nonnegativity.
Since $g(k+1)-g(k)=t_k$, subtracting successive prefix sums gives the
Poisson equation for $k<R$. The full mean-zero sum gives the terminal
equation
\[
 Rg(R-1)=R-\mu.
\]
In the actual chain the birth rate of $y$ is $(b/D)(y+1)$ and its death
rate is $y(x+1)/D\le y$. Increasing a birth rate multiplies a nonpositive
increment of $g$; decreasing a death rate multiplies a nonnegative
increment. This proves (6), including $y=R$, whose upward increment
is zero. At $y=R+1$ the downward successor is $R$, where $g$ is also zero;
this proves the last assertion.
\end{proof}

\section{A single cutoff function and its interfaces}

Let $f$ be zero on $[0,\gamma]$, linear with slope $\ell$ on
$[\gamma,\eta]$, and one on $[\eta,\infty)$. Define
\[
 J_X=f(b/D)g(y),\quad J_Y=f(a/D)g(x),\quad J=J_X+J_Y,
 \qquad V=W+ch_m(Z)+J.                                       \tag{7}
\]
Thus $0\le J\le2G$ and $V\ge W\ge\mathsf n$.

\begin{lemma}\label{lem:interface}
At every state,
\begin{align}
 \LL J&\le f(b/D)\LL g(y)+f(a/D)\LL g(x)
                  +\frac{\ell G(\lambda+4d)}D,               \tag{8}\\
 \LL J&\le G(\lambda+2d).                                    \tag{9}
\end{align}
\end{lemma}
\begin{proof}
For a transition $s\to s'$ the exact product identity is
\[
 \Delta J_X=f(b/D)\,[g(y')-g(y)]
              +g(y')\,[f(b'/D')-f(b/D)].                      \tag{10}
\]
There is an identical expression for $J_Y$.
A first download increases $D$ and does not increase either waiting-room
numerator. Both cutoff increments are nonpositive, so their interface
contributions in (10) are nonpositive, regardless of $a,b$.
An arrival changes only one cutoff, by at most $\ell/D$.

A departure has $D\ge2$ and decreases $D$ by one. For either numerator
$w$,
\[
 0\le f\!\left(\frac w{D-1}\right)-f\!\left(\frac wD\right)
       \le \frac{2\ell}D.                                    \tag{11}
\]
Indeed the difference is zero if $w/D\ge\eta$; otherwise
$w<\eta D$ and the Lipschitz bound is
$\ell w/[D(D-1)]\le \ell\eta/(D-1)\le2\ell/D$.
There are two cutoff terms, and the post-jump $g$ values are at most
$G$. Summing their interface bounds over the rates gives (8).
There are no departures at $D=1$.

For (9), a first download makes each whole product in (7)
nonincreasing: its cutoff decreases and its affected $g$ value also
decreases. An arrival increases $J$ by at most $G$, and a departure
increases it by at most $2G$. Sum over their total rates $\lambda,d$.
\end{proof}

\section{Negative drift in all states outside a finite set}

Choose a finite integer $N$ satisfying
\begin{align*}
 N&\ge R+2,\qquad N\ge3(R+1),\\
 (1-\eta)N-(1+\eta)R-\eta&\ge m,\\
 N&\ge\ell G[\lambda+4(2R+1)].
\end{align*}
All constants are finite and $1-\eta>0$.

If $\min(x,y)\ge R+1$, both corrections and every one-step successor
correction vanish. Also $d\ge\min(x,y)$: for $x\ge y$,
$D(d-y)=y(x-y)+x\ge0$. Equations (4)--(5) imply
\[
 \LL V\le\beta+c-q(R+1)<-2.                                  \tag{12}
\]

Next suppose $y\le R$ and $x\ge N$.
The correction involving $g(x)$ and all its successors vanish.
We have
\[
 d\le2R+1,\qquad \frac{x-y}D\ge\frac12,\qquad qd\ge y.          \tag{13}
\]
For the last inequality, $x\ge3(y+1)$ gives $x/D\ge3/4$, so
$qd\ge(21/32)(2y+1)\ge y$.
The remainder in (8) is at most one by the choice of $N$.
Lemma \ref{lem:poisson}, including the case $f=0$, gives
\[
 f(b/D)\LL g(y)\le f(b/D)(y-\mu).
\]
If $b/D\ge\eta$, the cutoff is one and
\[
 \LL V\le\beta-u-qd+c+y-\mu+1
             \le\beta+c-\mu+1<-3.                            \tag{14}
\]
If $b/D<\eta$, then
\[
 Z=a+x-b-y\ge(1-\eta)x-(1+\eta)y-\eta\ge m.
\]
The last term in (4) is therefore at most $-c/2$.
Since $f(b/D)(y-\mu)\le y$ and $qd\ge y$,
\[
 \LL V\le\beta-u-c/2+1=-3-u\le-3.                             \tag{15}
\]
This includes the entire cutoff ramp. Interchanging the two types
handles the opposite corner.

In every remaining state $x+y\le N+R$. Since $d\le x+y$, (4) and (9)
give
\[
 \LL V\le\beta+c+G[\lambda+2(N+R)]-u,\qquad
 u\ge\frac{a+b}{N+R+1}.
\]
With
\[
 B=(N+R+1)\{\beta+c+G[\lambda+2(N+R)]+1\},
\]
we conclude
\[
 \LL V\le-1\quad\text{outside}\quad
 F=\{x+y\le N+R,\ a+b\le B\}.                                \tag{16}
\]
The set $F$ is finite. This is one globally defined certificate,
rather than a choice between unaccounted local drift functions.

\section{The continuous-time recurrence conclusion}

We give the recurrence step to make its clock and integrability
requirements explicit. Arrivals form a rate-$\lambda$ Poisson process.
Every internal event decreases $W$ by one and every arrival increases
it by two. Thus the total number of jumps up to time $t$ is at most
$W(0)+3P_\lambda(t)$, proving nonexplosion. Every state has a
positive-rate finite path to zero, because the seed makes all required
downloads and departures possible. Conversely, arrivals followed by
prescribed first downloads build any finite state, proving
irreducibility.

Since $|Z|\le W$ and $h_m(z)\le |z|+m/2$, we have
\[
 W\le V\le (1+c)W+cm/2+2G.                                   \tag{17}
\]
Every jump of $V$ is bounded in absolute value by a finite constant
depending on $\lambda$. Localize Dynkin's identity at the first hit
$\tau_F$ of $F$, time $t$, and the first exit from a finite population
set. Equation (16) and $V\ge0$ bound the expected stopped time by
$V(s)$. Nonexplosion and Fatou remove the population cutoff; monotone
convergence as $t\to\infty$ yields
\[
                  \E_s\tau_F\le V(s),\qquad s\notin F.        \tag{18}
\]
No invariant law or stationary moment was used.

For completeness, select from every state of finite $F$ a finite
positive-rate path to zero. These finitely many paths have a common
time horizon $T$ and a positive lower success probability.
Starting in $F$, the bound
$W(T)\le W(0)+2P_\lambda(T)$ and (17) give a uniform finite bound on
$\E V(S_T)$. Equation (18) therefore bounds uniformly the expected
duration of a trial of length $T$ followed by return to $F$.
The strong Markov property and the uniform positive success probability
give finite expected hitting time of zero by repeated trials.
The mean holding time at zero is $1/\lambda$; each of its two exit
neighbors has finite expected hitting time back to zero.
The resulting regeneration cycle has finite mean. Its normalized
occupation measure is stationary, and irreducibility gives uniqueness.
This proves Theorem \ref{thm:main}.

\section{Credit, verification, and disclosure}

The model, its conjecture, the distinction between stochastic and
large-system stability, and the standard Foster principle are prior
work \cite{OWR,NRE09,NRE11,MT}. The finite geometric birth--death
Poisson equation and Lipschitz rounding estimate are standard tools.
The certificate above combines them with cutoffs whose first-download
interface terms have a useful sign.

Oguz, Anantharam and Norros \cite{OAN12} prove stability of their
Common Chunk protocol using samples of three from the entire
population and a rule protecting departure of rare chunks.
That protocol is different from (1). For example, with just the seed
and one peer possessing chunk one, its departure probability per
unit contact time is $7/8$, whereas (1) has departure rate $1/2$.
A fixed arrival rate also prevents identifying the chains by a
state-dependent clock change. Their final remaining rare-chunk
conjecture concerns a different algorithm and is not evidence that
the deterministic-first-chunk conjecture remained open.

An extensive bounded priority review identified no earlier exact all-load
recurrence theorem in the sources inspected. The complete 2012 Allerton
body was compared with this generator, including its operative proof.
The final 2011 publisher body, the final 2015 journal extension of the
2012 paper, and some later bodies remain uninspected. Their absence
from the checked evidence is not proof that they contain no relevant
result. The supplement records actual reading and access limits. No
worldwide-firstness or continuing-openness certificate is asserted.
The source target is the stated 2010 conjecture and the checked 2011
author-proof model.

The supporting materials preserve the original proof and exact
rational control programs, fresh independent boundary and
counterexample controls, expected outputs, source-version notes and
file hashes. They check finite states and identities; the all-state,
all-real-load theorem is proved analytically above. These controls
are not a proof-assistant certificate. No stationary moment or
performance bound is inferred from them.

AI tools were used extensively in the research, proof development,
literature comparison, computational reproduction, manuscript
preparation and adversarial auditing. This is an unrefereed preprint
draft and has not undergone formal independent human peer review.
Internal AI reviews do not constitute such human peer review.

\begin{thebibliography}{9}
\bibitem{OWR}
I.~Norros, On the stability of population processes
(joint work with H.~Reittu), in \emph{Mathematical Challenges in
Stochastic Networks}, Oberwolfach Reports 48/2010, pp.~2782--2783.
\href{https://doi.org/10.4171/OWR/2010/48}{doi:10.4171/OWR/2010/48}.
\bibitem{NRE09}
I.~Norros, H.~Reittu and T.~Eirola,
\emph{On the stability of two-chunk file-sharing systems},
arXiv:0910.5577v1 (2009).
\url{https://arxiv.org/abs/0910.5577v1}.
\bibitem{NRE11}
I.~Norros, H.~Reittu and T.~Eirola,
On the stability of two-chunk file-sharing systems,
\emph{Queueing Systems} \textbf{67} (2011), 183--206.
\href{https://doi.org/10.1007/s11134-011-9209-2}
{doi:10.1007/s11134-011-9209-2}.
Model and conjecture checked in the January 5, 2011 author proof,
\url{https://hal.science/hal-00781341}.
\bibitem{MT}
S.~P.~Meyn and R.~L.~Tweedie,
Stability of Markovian processes III: Foster--Lyapunov criteria for
continuous-time processes,
\emph{Advances in Applied Probability} \textbf{25} (1993), 518--548.
\href{https://doi.org/10.2307/1427522}{doi:10.2307/1427522}.
\bibitem{OAN12}
B.~Oguz, V.~Anantharam and I.~Norros,
Stable, distributed P2P protocols based on random peer sampling,
in \emph{50th Annual Allerton Conference on Communication, Control,
and Computing} (2012), pp.~915--919.
\end{thebibliography}
\end{document}
